% ==============================================================================
%  SCD · Kit del parcial del Tema 2.2: algoritmos de exclusión mutua
%  Examen: 26 de noviembre de 2026 (convocado el 19, se cambió de fecha)
%  Fuentes: diapositivas «Tema 2. Sincronización en memoria compartida» (22/10/2026)
%  y Relación de problemas Tema 2-2 (problemas 54–63), Prof. M. I. Capel (UGR).
%  Compilar: pdflatex kit_parcial_t22.tex (dos pasadas)
% ==============================================================================
\documentclass[11pt,a4paper,oneside]{report}

% ==============================================================================
% CONFIGURACIÓN Y ESTILO — RELACIONES DE PROBLEMAS RESUELTAS DE SCD
% ==============================================================================

\usepackage[utf8]{inputenc}
\usepackage[T1]{fontenc}
\usepackage[spanish,es-tabla]{babel}
\usepackage{lmodern}
\usepackage{microtype}
\usepackage[table,xcdraw]{xcolor}
\usepackage{amsmath,amssymb,amsfonts,amsthm}
\usepackage{graphicx}
\usepackage{booktabs}
\usepackage{array}
\usepackage{multirow}
\usepackage{enumitem}
\usepackage{fancyhdr}
\usepackage{titlesec}
\usepackage{listings}
\usepackage{tcolorbox}
\tcbuselibrary{skins,breakable}
\usepackage{tikz}
\usetikzlibrary{arrows.meta,positioning,shapes,babel}

% Configuración de márgenes
\usepackage[a4paper,top=2.5cm,bottom=2.5cm,left=2.6cm,right=2.6cm,headheight=24pt]{geometry}

% Hipervínculos
\usepackage{hyperref}
\hypersetup{
    colorlinks=true,
    linkcolor=scdDarkBlue,
    citecolor=scdTeal,
    urlcolor=scdAccent,
    pdfauthor={repasaYA - UGR},
    pdftitle={SCD - Relaciones de Problemas Resueltas},
    pdfsubject={Sistemas Concurrentes y Distribuidos},
    pdfkeywords={SCD, Ejercicios, Hoare, Concurrencia, Semáforos, Monitores, MPI, Tiempo Real}
}

% Paleta de colores
\definecolor{scdDarkBlue}{RGB}{20, 35, 60}
\definecolor{scdPrimary}{RGB}{35, 75, 135}
\definecolor{scdAccent}{RGB}{108, 99, 255}
\definecolor{scdLightAccent}{RGB}{240, 242, 255}
\definecolor{scdGreen}{RGB}{28, 140, 72}
\definecolor{scdLightGreen}{RGB}{242, 253, 245}
\definecolor{scdOrange}{RGB}{217, 119, 6}
\definecolor{scdLightOrange}{RGB}{255, 251, 235}
\definecolor{scdRed}{RGB}{185, 28, 28}
\definecolor{scdGray}{RGB}{246, 248, 251}
\definecolor{scdBorderGray}{RGB}{210, 218, 230}
\definecolor{scdTeal}{RGB}{15, 118, 110}
\definecolor{codeBg}{RGB}{248, 249, 252}

% Encabezados y Pies
\pagestyle{fancy}
\fancyhf{}
\fancyhead[LE,RO]{\footnotesize\textsf{\thepage}}
\fancyhead[LO]{\footnotesize\textsf{\rightmark}}
\fancyhead[RE]{\footnotesize\textsf{\leftmark}}
\renewcommand{\headrulewidth}{0.4pt}
\renewcommand{\footrulewidth}{0pt}

% Formato de Capítulos (Relaciones)
\titleformat{\chapter}[display]
  {\normalfont\bfseries\color{scdDarkBlue}\sffamily}
  {\Large\textsf{Relación de Problemas \thechapter}}{1ex}
  {\titlerule[1.5pt]\vspace{1ex}\huge}
  [\vspace{1ex}\titlerule]

\titleformat{\section}
  {\normalfont\Large\bfseries\color{scdPrimary}\sffamily}
  {\thesection}{1em}{}

\titleformat{\subsection}
  {\normalfont\large\bfseries\color{scdDarkBlue}\sffamily}
  {\thesubsection}{1em}{}

% Entorno: Metadatos y Fuentes
\newcommand{\cuadrofuentes}[1]{%
  \begin{tcolorbox}[
    colback=scdLightAccent,
    colframe=scdAccent,
    boxrule=1pt,
    arc=4pt,
    left=12pt,
    right=12pt,
    top=8pt,
    bottom=8pt,
    breakable
  ]
    \small
    \noindent\textbf{\color{scdAccent}\large\textsf{Metadatos de la Relación y Fuentes:}}\\[4pt]
    #1
  \end{tcolorbox}
  \vspace{10pt}
}

% Entorno: Enunciado del Ejercicio
\newtcolorbox{ejercicio}[1][]{
  enhanced,
  colback=scdGray,
  colframe=scdPrimary,
  boxrule=1pt,
  leftrule=4pt,
  arc=3pt,
  title=\textbf{\textsf{Ejercicio #1}},
  coltitle=white,
  colbacktitle=scdPrimary,
  attach boxed title to top left={yshift=-2mm, xshift=4mm},
  boxed title style={arc=2pt, boxrule=0pt},
  top=10pt,
  bottom=8pt,
  left=10pt,
  right=10pt,
  breakable
}

% Entorno: Solución Explicada Paso a Paso
\newtcolorbox{solucion}[1][]{
  enhanced,
  colback=scdLightGreen,
  colframe=scdGreen,
  boxrule=0.8pt,
  leftrule=3.5pt,
  arc=3pt,
  title=\textbf{\textsf{Solución Detallada\ifstrempty{#1}{}{: #1}}},
  coltitle=white,
  colbacktitle=scdGreen,
  attach boxed title to top left={yshift=-2mm, xshift=4mm},
  boxed title style={arc=2pt, boxrule=0pt},
  top=10pt,
  bottom=8pt,
  left=10pt,
  right=10pt,
  breakable
}

% Entorno: Aviso / Ojo Examen
\newtcolorbox{aviso}[1][]{
  enhanced,
  colback=scdLightOrange,
  colframe=scdOrange,
  boxrule=0.8pt,
  leftrule=3.5pt,
  arc=3pt,
  title=\textbf{\textsf{Clave de Examen\ifstrempty{#1}{}{: #1}}},
  coltitle=white,
  colbacktitle=scdOrange,
  attach boxed title to top left={yshift=-2mm, xshift=4mm},
  boxed title style={arc=2pt, boxrule=0pt},
  top=10pt,
  bottom=8pt,
  left=10pt,
  right=10pt,
  breakable
}

% Listings para código
\lstset{
  backgroundcolor=\color{codeBg},
  basicstyle=\ttfamily\footnotesize\color{black},
  keywordstyle=\bfseries\color{scdPrimary},
  commentstyle=\itshape\color{gray},
  stringstyle=\color{scdGreen},
  numbers=left,
  numberstyle=\tiny\color{gray},
  numbersep=8pt,
  stepnumber=1,
  breaklines=true,
  breakatwhitespace=false,
  showstringspaces=false,
  frame=none,
  tabsize=2,
  captionpos=b,
  inputencoding=utf8,
  extendedchars=true,
  literate=
    {á}{{\'a}}1 {é}{{\'e}}1 {í}{{\'i}}1 {ó}{{\'o}}1 {ú}{{\'u}}1
    {Á}{{\'A}}1 {É}{{\'E}}1 {Í}{{\'I}}1 {Ó}{{\'O}}1 {Ú}{{\'U}}1
    {ñ}{{\~n}}1 {Ñ}{{\~N}}1 {¿}{{?`}}1 {¡}{{!`}}1
    {«}{{\guillemotleft}}1 {»}{{\guillemotright}}1
}


\renewcommand{\thesection}{\arabic{section}}
\setcounter{secnumdepth}{2}
\pagestyle{fancy}
\fancyhf{}
\fancyhead[R]{\footnotesize\textsf{\thepage}}
\fancyhead[L]{\footnotesize\textsf{SCD · Parcial Tema 2.2 · Algoritmos de exclusión mutua}}

\lstset{language={},numbers=none,basicstyle=\ttfamily\small\color{black}}

\hypersetup{
  pdftitle={SCD - Kit del parcial del Tema 2.2: algoritmos de exclusion mutua},
  pdfauthor={repasaYA - UGR},
  pdfsubject={Sistemas Concurrentes y Distribuidos}
}

\newcommand{\si}{\textcolor{scdGreen}{\textbf{Sí}}}
\newcommand{\no}{\textcolor{scdRed}{\textbf{No}}}

\begin{document}

% ------------------------------------------------------------------------------
\begin{titlepage}
  \centering
  \vspace*{1.5cm}
  {\Huge\bfseries\color{scdDarkBlue}\textsf{Sistemas Concurrentes y Distribuidos}}\\[0.6cm]
  {\LARGE\bfseries\color{scdPrimary}\textsf{Kit del parcial del Tema 2.2}}\\[0.3cm]
  {\Large\textsf{Algoritmos de exclusión mutua}}\\[1.2cm]
  \begin{tcolorbox}[colback=scdLightAccent,colframe=scdAccent,arc=4pt,width=0.86\textwidth]
    \centering\large
    \textbf{Jueves 26 de noviembre de 2026} \quad·\quad 17:30 \quad·\quad aula 0.2\\[2pt]
    {\normalsize (la convocatoria decía 19 de noviembre; se cambió al 26)}\\[6pt]
    \textbf{60 minutos} \quad·\quad sin apuntes ni dispositivos\\[6pt]
    35 preguntas test (5,25 puntos) \;+\; 3 ejercicios breves obligatorios (4,75 puntos)\\[2pt]
    Acierto $+0{,}15$ \quad Fallo $-0{,}075$ \quad En blanco $0$ \quad (la parte test no baja de 0)
  \end{tcolorbox}
  \vspace{1cm}
  \begin{minipage}{0.86\textwidth}\small
  \textbf{Qué entra.} Lo que dicen las diapositivas del 22/10/2026: el problema de la exclusión mutua y las
  condiciones de Dijkstra, el método de refinamiento sucesivo hasta Dekker, los algoritmos de Dijkstra, Knuth y
  Peterson para $N$ procesos, y los algoritmos distribuidos de Ricart-Agrawala, Suzuki-Kasami y Misra. Más la
  Relación 2-2 (problemas 54–63): Hyman, tickets, anillo, \texttt{swap} y cerrojo de capacidad $k$.\\[6pt]
  \textbf{Cómo usar este kit.} Sección 1 para la táctica; secciones 2–7 para la teoría tal como la pregunta el
  test; sección 8 con las plantillas de los ejercicios breves y 9 ejercicios resueltos; y la tabla resumen final
  para el último repaso. Las soluciones completas de la relación están en \emph{Relaciones de problemas
  resueltas} (problemas 54–63) y el banco de 82 preguntas en el \emph{Simulacro parcial T2.2} de repasaYA.
  \end{minipage}
  \vfill
  {\large\textsf{repasaYA · Universidad de Granada}}
\end{titlepage}

\tableofcontents
\clearpage

% ==============================================================================
\section{Táctica para los 60 minutos}

\subsection{Reparto del tiempo}
\begin{center}
\begin{tabular}{lll}
\toprule
\textbf{Bloque} & \textbf{Tiempo} & \textbf{Cómo} \\ \midrule
Test (35 preguntas) & 25 min & unos 40 segundos por pregunta; dos pasadas \\
3 ejercicios breves & 30 min & 10 minutos cada uno, empezando por el que veas más claro \\
Repaso & 5 min & revisa que cada ejercicio tenga conclusión explícita \\
\bottomrule
\end{tabular}
\end{center}
Primera pasada del test: contesta solo lo que sabes seguro y marca las dudas. Segunda pasada: aplica la regla
de abajo. No dejes el test para el final: vale más que los ejercicios.

\subsection{¿Contestar o dejar en blanco?}
Un fallo resta exactamente \textbf{medio acierto} ($0{,}075 = 0{,}15/2$). La esperanza de contestar al azar entre
las opciones que te quedan es:
\[
E = \tfrac{1}{m}\cdot 0{,}15 - \tfrac{m-1}{m}\cdot 0{,}075 \qquad (m = \text{opciones que no has descartado})
\]
\begin{center}
\begin{tabular}{cccc}
\toprule
\textbf{Te quedan} & \textbf{Esperanza} & \textbf{Decisión} \\ \midrule
2 opciones & $+0{,}0375$ & \textbf{contesta} \\
3 opciones & $0$ & indiferente: contesta solo si una te parece más probable \\
4 opciones & $-0{,}019$ & en blanco \\
5 opciones & $-0{,}030$ & en blanco \\
\bottomrule
\end{tabular}
\end{center}
\begin{aviso}[regla corta]
Si dudas entre dos, contesta. Si no has podido descartar nada, en blanco.
\end{aviso}

\subsection{Los ejercicios breves}
Valen 4,75 puntos entre tres (unos 1,6 cada uno). Lo que se puntúa es \textbf{justificar}: la relación lo dice
explícitamente, «un contraejemplo debe incluir una secuencia de ejecución suficientemente detallada». Casi todos
son de uno de estos tipos (plantillas en la sección~\ref{sec:plantillas}):
\begin{enumerate}[leftmargin=*]
  \item \textbf{Encontrar un fallo}: una tabla de intercalación que lleva a dos procesos a la SC, a un
        interbloqueo o a un livelock.
  \item \textbf{Demostrar una propiedad}: exclusión mutua por reducción al absurdo; progreso e inanición
        indicando la hipótesis de planificación.
  \item \textbf{Analizar una modificación o diseñar}: cambiar una guarda, intercambiar dos asignaciones,
        protocolo con \texttt{swap}, cerrojo de capacidad $k$, o seguir un escenario de Ricart-Agrawala.
\end{enumerate}

% ==============================================================================
\section{El problema y las condiciones de Dijkstra}

\subsection{Historia (cae en el test)}
\begin{center}\small
\begin{tabular}{lp{11.5cm}}
\toprule
\textbf{Año} & \textbf{Aportación} \\ \midrule
1962 & \textbf{Dekker} plantea el problema para multiprocesadores: acceso mutuamente excluyente a una SC, sin interbloqueo. \\
1965 & \textbf{Dijkstra}: solución segura y libre de interbloqueo, pero \textbf{puede producir inanición}. \\
1966 & \textbf{Knuth}: \textbf{sin inanición}, retraso limitado, pero \textbf{no FIFO}. \\
1974 & \textbf{Lamport}: los procesos pueden detenerse en el protocolo, se solapan lecturas y escrituras, y hay retraso \textbf{FIFO} para los que ya esperan. \\
1981 & \textbf{Peterson}: equitativa para 2 y $n$ procesos, \textbf{retraso cuadrático}, la más simple. \\
1983 & Algoritmos \textbf{distribuidos} para multicomputadores: Ricart-Agrawala, Suzuki-Kasami. \\
\bottomrule
\end{tabular}
\end{center}

\subsection{Las cuatro propiedades y su tipo}
\begin{center}\small
\begin{tabular}{p{5cm}p{3.4cm}p{6cm}}
\toprule
\textbf{Propiedad} & \textbf{Tipo} & \textbf{Hipótesis de planificación} \\ \midrule
Exclusión mutua & seguridad & ninguna (vale para cualquier intercalación) \\
Ausencia de interbloqueo (progreso, alcanzabilidad) & vivacidad (global) & progreso finito \\
Ausencia de inanición & vivacidad (individual) & planificación (débilmente) justa \\
Espera acotada (equidad) & vivacidad + cota & cota de adelantamientos, no de tiempo \\
\bottomrule
\end{tabular}
\end{center}
\textbf{EM + ausencia de interbloqueo no implica ausencia de inanición}: el sistema progresa (alguien entra),
pero un proceso concreto puede ser adelantado siempre. Ejemplo de clase: el algoritmo de Dijkstra para $N$.

\textbf{Livelock (interbloqueo activo)}: los procesos siguen ejecutando instrucciones pero ninguno avanza hacia
la SC (4ª etapa del refinamiento).

\subsection{Condiciones de Dijkstra (solución parcialmente correcta)}
\begin{enumerate}[leftmargin=*]
  \item No hacer ninguna suposición sobre las instrucciones o el número de procesos del multiprocesador.
  \item Ni sobre la velocidad de ejecución de los procesos, excepto que no es cero (\textbf{progreso finito}).
  \item Un proceso que ejecuta código \textbf{fuera} de la SC no puede impedir a los otros entrar en ella.
  \item La SC siempre será alcanzada por alguno de los procesos que esperan entrar (\textbf{alcanzabilidad}).
\end{enumerate}
Un algoritmo que solo cumple estas cuatro condiciones es \textbf{parcialmente correcto}: tiene EM y
alcanzabilidad, pero \textbf{puede tener inanición}. Lo que añaden Knuth, Peterson y Lamport es la equidad.

\textbf{Espera activa}: los procesos iteran en un bucle vacío hasta que entrar es seguro. Funciona, pero
consume CPU; solo es aceptable si hay pocos procesos.

% ==============================================================================
\section{Refinamiento sucesivo hasta Dekker (2 procesos)}
En las diapositivas la clave $c_i = 0$ significa «quiero entrar o estoy dentro» y $c_i = 1$ «estoy fuera».

\begin{center}\small
\begin{tabular}{p{1.3cm}p{5.6cm}p{6.6cm}}
\toprule
\textbf{Etapa} & \textbf{Idea} & \textbf{Qué falla} \\ \midrule
1ª & Variable \texttt{turno} que alterna (corrutinas). & EM sí, pero \textbf{viola la 3ª condición}: si uno no quiere entrar, el otro no puede (alternancia forzosa). \\
2ª & Espero mientras $c_j=0$; \textbf{después} pongo $c_i:=0$. & \textbf{Falla la EM}: la salida de la espera y el cambio de clave no son atómicos. \\
3ª & Pongo $c_i:=0$ \textbf{antes} de mirar $c_j$. & Segura, pero \textbf{interbloqueo}: los dos ponen su clave a 0 y esperan para siempre. \\
4ª & Si $c_j=0$, cedo: $c_i:=1$, espero, $c_i:=0$ y reintento. & EM sí; sin bloqueo permanente, pero \textbf{livelock/inanición} si van al mismo ritmo. \\
5ª & \textbf{Dekker}: claves + \texttt{turno}. Solo cede el que no tiene el turno. & \textbf{Correcto} para 2 procesos: EM, sin interbloqueo, sin inanición. \\
\bottomrule
\end{tabular}
\end{center}

\subsection{Algoritmo de Dekker (proceso $P_1$; $P_2$ es simétrico)}
\begin{lstlisting}
c1 := 0;                       -- quiero entrar
while c2 = 0 do                -- el otro tambien quiere
  if turno = 2 then begin      -- y el turno es suyo: cedo
    c1 := 1;
    while turno = 2 do nothing; enddo;
    c1 := 0;
  end endif;
enddo;
<<seccion critica>>
turno := 2;                    -- cedo el turno al salir
c1 := 1;
\end{lstlisting}
\begin{itemize}[leftmargin=*]
  \item \textbf{EM}: $P_i$ solo entra si ve $c_j = 1$, y lo comprueba \emph{después} de asignar $c_i = 0$.
        Al entrar se cumple $c_j = 1 \land c_i = 0$; si los dos estuvieran dentro, cada uno habría visto la clave
        del otro a 1 después de que el otro la pusiera a 0. Imposible. (El turno no interviene en la EM.)
  \item \textbf{Alcanzabilidad}: si los dos quieren entrar y $\texttt{turno} = i$, $P_i$ espera a que $P_j$ ceda su
        clave y entra; $P_j$ pone su clave a 1 y queda en espera activa.
  \item \textbf{Sin inanición}: al salir se cede el turno, así que el que salió no puede adelantar
        indefinidamente al que espera (con planificación justa).
\end{itemize}

% ==============================================================================
\section{Dijkstra y Knuth para $N$ procesos}

\subsection{Algoritmo de Dijkstra}
\texttt{c: array[0..n-1] of (pasivo, solicitando, en\_SC); turno: 0..n-1}
\begin{lstlisting}
repeat
  repeat
E2: c[i] := solicitando;
    while turno <> i do                      -- 1a barrera
E3:   if c[turno] = pasivo then turno := i endif;
    enddo;
E4: c[i] := en_SC;
    j := 0;                                  -- 2a barrera
    while (j < n) and (j = i or c[j] <> en_SC) do j := j+1; enddo;
  until j >= n;
  <<seccion critica>>
E1: c[i] := pasivo;
until false
\end{lstlisting}
\begin{itemize}[leftmargin=*]
  \item \textbf{1ª barrera}: detiene a los procesos si el que posee el turno no está pasivo. La comprobación y el
        cambio de turno \textbf{no son atómicos}: pueden pasar varios a la vez.
  \item \textbf{2ª barrera}: da la \textbf{seguridad}. Se marca \texttt{en\_SC} y solo se entra si nadie más lo
        está; si no, se vuelve a empezar.
  \item \textbf{Alcanzabilidad}: \texttt{turno} conserva el valor del último que lo asignó. Entre los que llegan a
        \texttt{en\_SC}, entra el que tiene el turno y los demás vuelven a la 1ª barrera.
  \item \textbf{Inanición: sí puede haberla}. Un proceso puede quedarse siempre en la 1ª barrera mientras otros
        toman el turno (traza de la diapositiva 2.13).
\end{itemize}

\begin{aviso}[problema 63: cambiar la guarda a \texttt{c[turno] <> en\_SC}]
Con la nueva guarda se puede quitar el turno a un proceso que está \emph{solicitando}. Dos procesos pueden
quitárselo uno a otro indefinidamente: \textbf{se pierde la alcanzabilidad} (progreso). La \textbf{EM se mantiene}
porque la 2ª barrera no cambia. Ojo: en un banco de preguntas de otros años la clave marcaba «pierde corrección
(seguridad)». Si en el test hay una opción de progreso/vivacidad, elige esa; si la única que habla de pérdida
es «pierde corrección», es la que más se acerca.
\end{aviso}

\subsection{Algoritmo de Knuth}
\begin{lstlisting}
repeat
  repeat
E0: c[i] := solicitando;
    j := turno;                              -- variable local
E1: while j <> i do                          -- 1a barrera
      if c[j] <> pasivo then j := turno
      else j := (j-1) mod n endif;
    enddo;
E2: c[i] := en_SC;
    k := 0;                                  -- 2a barrera (igual que Dijkstra)
    while (k < n) and (k = i or c[k] <> en_SC) do k := k+1; enddo;
  until k >= n;
E3: turno := i;
  <<seccion critica>>
  turno := (i-1) mod n;                      -- cede el turno al predecesor
E4: c[i] := pasivo;
until false
\end{lstlisting}
\begin{itemize}[leftmargin=*]
  \item \textbf{Novedad}: el \textbf{turno circular}. El que espera recorre el anillo desde el turno hacia atrás y
        el que sale cede el turno a su predecesor.
  \item \textbf{Cumple EM, alcanzabilidad y ausencia de inanición} (retraso limitado), \textbf{no FIFO}.
  \item La imposibilidad de inanición exige un \textbf{número finito de procesos}: con infinitos, uno podría ser
        adelantado siempre.
  \item Dato clásico (no está en las diapositivas): la espera máxima de Knuth es exponencial, $2^{n-1}-1$
        turnos. Por eso Peterson ($n(n-1)/2$, cuadrático) supone una mejora.
\end{itemize}

% ==============================================================================
\section{Peterson para 2 y para $N$ procesos}

\subsection{Dos procesos}
Versión de las diapositivas (el que escribe \texttt{turno} el último es el que espera):
\begin{lstlisting}
solicitado[i] := true;
turno := i;
while (solicitado[j] and turno = i) do nothing; enddo;
<<seccion critica>>
solicitado[i] := false;
\end{lstlisting}
En la relación (problema 57) aparece como \texttt{turno := 1-i} y se espera mientras
\texttt{quiere[1-i] and turno = 1-i} («cedo el turno»). Es el mismo algoritmo: \textbf{espera el último que
escribe \texttt{turno}}.
\begin{itemize}[leftmargin=*]
  \item \textbf{EM} (por reducción al absurdo): si los dos están dentro, los dos indicadores son ciertos, así que
        cada uno tuvo que ver el turno a su favor. Pero \texttt{turno} tiene un único valor, el del último que
        lo escribió, y ese último no pudo pasar. Contradicción.
  \item \textbf{Progreso}: si solo uno quiere entrar, entra sin esperar (el otro indicador es falso). No hay
        alternancia forzosa.
  \item \textbf{Sin inanición y espera acotada}: si el rival sale y vuelve a pedir, al escribir el turno se
        convierte en el último y cede. \textbf{A lo sumo un adelantamiento}.
  \item \textbf{Intercambiar las dos primeras asignaciones rompe la EM} (ejercicio resuelto 2).
  \item \textbf{Memoria débil}: la demostración supone consistencia secuencial. En hardware real la lectura del
        indicador del otro puede adelantarse a la escritura del propio; hay que usar atómicos con orden
        secuencialmente consistente o barreras.
\end{itemize}

\subsection{$N$ procesos: el filtro por etapas}
\texttt{c: array[0..n-1] of -1..n-2}: etapa de cada proceso ($-1$ = fuera).\\
\texttt{turno: array[0..n-2] of 0..n-1}: último proceso en llegar a cada etapa.
\begin{lstlisting}
for j := 0 to n-2 do
  c[i] := j;
  turno[j] := i;                       -- el ultimo en llegar se queda con el turno
  while (exists k <> i: c[k] >= j) and turno[j] = i do nothing; enddo;
enddo;
c[i] := n-1;                           -- metainstruccion: esta en la SC
<<seccion critica>>
c[i] := -1;
\end{lstlisting}
Un proceso deja de esperar en la etapa $j$ si (a) va por delante de todos o (b) llega otro después a su etapa.

\begin{center}\small
\begin{tabular}{lp{11.8cm}}
\toprule
\textbf{Lema} & \textbf{Enunciado} \\ \midrule
Def. & $P_i$ \textbf{precede} a $P_j$ si $c[i] > c[j]$. \\
L1 & Un proceso que precede a todos los demás puede avanzar al menos una etapa. \\
L2 & Quien pasa de la etapa $j$ a la $j+1$ o precede a todos o no estaba solo en la etapa $j$. \\
L3 & Si hay al menos dos procesos en la etapa $j$, hay al menos uno en cada etapa anterior (inducción). \\
L4 & En la etapa $j$ hay como máximo $n - j$ procesos ($0 \le j \le n-2$). A la etapa $n-2$ llegan como máximo 2. \\
\bottomrule
\end{tabular}
\end{center}
\begin{itemize}[leftmargin=*]
  \item \textbf{EM}: a la etapa $n-2$ llegan como máximo dos procesos (L4) y, por L2, solo uno puede pasar a la SC.
  \item \textbf{Alcanzabilidad}: por reducción al absurdo. Si nadie avanzara, el más adelantado o precede a todos
        (contradice L1) o comparte etapa (contradice L2).
  \item \textbf{Equidad}: espera máxima $r(n) = n - 1 + r(n-1) = \dfrac{n(n-1)}{2}$ turnos.
        (En un banco de otros años aparece «$N-1$»; según estas diapositivas es $n(n-1)/2$.)
  \item ¿Dos procesos en $n-2$ y ninguno en $n-3$? \textbf{No}, por L3.
\end{itemize}

% ==============================================================================
\section{Otros algoritmos de la relación}
\begin{itemize}[leftmargin=*]
  \item \textbf{Hyman (problema 56)}: \textbf{incorrecto}, viola la EM (ejercicio resuelto 1). Fallo de diseño: se
        comprueba la clave del otro y \emph{después}, sin volver a comprobar, se toma el turno; y quien ya tiene el
        turno entra sin mirar la clave del otro.
  \item \textbf{Dos indicadores sin turno (problema 55)}: EM \si{} (cada uno pone su indicador antes de leer el del
        otro); interbloqueo activo (livelock) \si{} posible si van a la par; inanición \si{} posible. El turno de
        Peterson es lo que rompe la simetría.
  \item \textbf{Tickets para 2 (problema 61)}: \texttt{n[i]:=1} avisa de que se está calculando el ticket; sin ella
        se pierde la EM. El desempate por identificador evita que los dos esperen con el mismo ticket. Ausencia de
        inanición por tickets crecientes y orden lexicográfico (ticket, id).
  \item \textbf{Panadería de Lamport}: número = 1 + máximo de los que ve; se atiende por (número, id); FIFO;
        los números crecen sin límite.
  \item \textbf{Anillo \texttt{s[i]} (problema 60)}: un testigo implícito que rota. EM \si; pero si un proceso se
        queda fuera de la SC, el testigo se para: \textbf{alternancia forzosa}, viola el progreso (3ª condición).
        La equidad solo acota adelantamientos entre los que compiten; la alternancia obliga aunque no quieras.
\end{itemize}

\section{Soporte hardware}
\begin{itemize}[leftmargin=*]
  \item \textbf{Deshabilitar interrupciones}: vale en un monoprocesador, no en un multiprocesador.
  \item \textbf{TestAndSet(x)}: devuelve el valor antiguo y pone \texttt{x := true}, atómicamente.
  \item \textbf{swap(x, y)}: intercambia atómicamente. Da EM, \textbf{no equidad}: un proceso puede perder la
        carrera siempre.
\end{itemize}

% ==============================================================================
\section{Algoritmos distribuidos}
Sin memoria compartida: se sincroniza con paso de mensajes. Se exige red \textbf{completamente conectada},
\textbf{sin errores}, con \textbf{retraso variable pero acotado}, y se admiten \textbf{desordenamientos} en la
entrega (no se exige FIFO).

\subsection{Ricart-Agrawala (permisos)}
\begin{itemize}[leftmargin=*]
  \item \textbf{Pedir}: \texttt{ns := mns + 1}; enviar \texttt{(pet, ns, i)} a los otros $n-1$; esperar
        $n-1$ respuestas (\texttt{num\_repesperadas}).
  \item \textbf{Al recibir \texttt{(pet, k, j)}}: \texttt{mns := max(ns, k)} (reloj lógico);
        \texttt{prioridad := solicitaSC and}\\ \texttt{(k > ns or (k = ns and i < j))}. Si tengo prioridad,
        \textbf{retraso} la respuesta (\texttt{repretrasadas[j] := true}); si no, contesto ya.
  \item \textbf{Al salir}: \texttt{solicitaSC := false} y respondo a todos los retrasados.
  \item \textbf{Mensajes por entrada}: $2(n-1)$.
  \item \textbf{EM}: si los dos entraran, cada uno habría contestado al otro después de generar su número, lo que
        es imposible: el de menor (\texttt{ns}, id) retrasa su respuesta.
  \item \textbf{Alcanzabilidad}: orden total de las peticiones (\texttt{ns}, id).
\end{itemize}

\subsection{Suzuki-Kasami (testigo)}
\begin{itemize}[leftmargin=*]
  \item Cada proceso: \texttt{token\_presente}, \texttt{en\_SC}, \texttt{peticion[1..N]}. El testigo lleva
        \texttt{token[1..N]} (número de la última petición atendida de cada uno).
  \item Si no tengo el testigo: \texttt{ns := ns+1}; difundo \texttt{(pet, ns, i)}; espero el testigo.
  \item Quien tiene el testigo y no está en la SC se lo envía a $j$ si \texttt{peticion[j] > token[j]}, buscando en
        orden circular $i+1, \dots, n, 1, \dots, i-1$ (\textbf{equidad}).
  \item \textbf{Mensajes}: $n$ ($n-1$ peticiones + 1 testigo); 0 si ya tengo el testigo.
  \item \textbf{EM}: globalmente hay exactamente un \texttt{token\_presente = true} (o el testigo está en tránsito).
  \item Pregunta abierta: ¿y si se pierden mensajes? Si se cambian los contadores por booleanos, se rompe la EM o la
        vivacidad (no se distinguen peticiones nuevas de atendidas).
\end{itemize}

\subsection{Misra (regeneración del testigo)}
\begin{itemize}[leftmargin=*]
  \item Dos testigos, \textbf{ping} y \textbf{pong}, recorren el anillo y cada uno vigila al otro.
  \item Inicialmente los dos en el mismo nodo con \texttt{num\_ping = 1}, \texttt{num\_pong = -1}.
        \textbf{Invariante}: \texttt{num\_ping + num\_pong = 0}.
  \item Cada proceso guarda en \texttt{m\_i} el número del último testigo que pasó por él. Si llega
        \texttt{ping} con \texttt{num\_ping = m\_i}, el pong no ha pasado desde entonces: \textbf{se perdió el
        pong} y se regenera (y simétricamente con pong).
  \item Cuando se encuentran, se incrementa el contador (número de colisiones) y se mantiene el invariante.
\end{itemize}

% ==============================================================================
\section{Plantillas y ejercicios breves resueltos}\label{sec:plantillas}

\subsection{Plantilla 1: contraejemplo (traza)}
\begin{enumerate}[leftmargin=*]
  \item Escribe el estado inicial de \emph{todas} las variables compartidas.
  \item Tabla: \textbf{paso · proceso · instrucción · variables tras el paso · comentario}. Una instrucción
        atómica por fila.
  \item Marca el punto clave («el planificador suspende a P0 aquí»).
  \item Termina con una frase: «\emph{Estado final: P0 y P1 en la SC. Se viola la EM.}» o «\emph{El estado se
        repite: livelock}».
\end{enumerate}

\subsection{Plantilla 2: demostrar la exclusión mutua}
\begin{enumerate}[leftmargin=*]
  \item Supón que $P_i$ y $P_j$ están a la vez en la SC.
  \item Escribe qué tuvo que ver cada uno para salir de su bucle de espera, y qué valen las variables que ninguno
        ha cambiado desde entonces (los indicadores siguen a cierto hasta salir).
  \item Fíjate en la \textbf{última escritura} de la variable de desempate (\texttt{turno}) o en el orden
        «escribo lo mío y después leo lo del otro».
  \item Llega a la contradicción y concluye: $N_{SC} \le 1$.
\end{enumerate}

\subsection{Plantilla 3: progreso e inanición}
Di siempre la \textbf{hipótesis de planificación}: «suponiendo planificación débilmente justa…». Para progreso:
si solo uno quiere entrar, ¿entra sin esperar? Si compiten, ¿quién pasa? Para inanición: ¿puede el rival
adelantarle una y otra vez? ¿Cuántas veces como máximo (espera acotada)?

\begin{ejercicio}[1: Hyman viola la exclusión mutua]
\texttt{c[0]=c[1]=1}, \texttt{turno=1}. $P_i$: \texttt{c[i]:=0; while turno<>i do (while c[j]=0 do skip; turno:=i);
SC; c[i]:=1}. Dar una intercalación que meta a los dos en la SC.
\end{ejercicio}
\begin{solucion}
\begin{center}\small
\begin{tabular}{clllc}
\toprule
\textbf{Paso} & \textbf{Proc.} & \textbf{Instrucción} & \textbf{c[0] c[1] turno} & \textbf{Comentario} \\ \midrule
1 & P0 & \texttt{c[0] := 0} & 0 \; 1 \; 1 & \\
2 & P0 & \texttt{turno <> 0}: cierto & 0 \; 1 \; 1 & entra al bucle \\
3 & P0 & \texttt{c[1] = 0}: falso & 0 \; 1 \; 1 & no espera \\
-- & & \multicolumn{3}{l}{\emph{el planificador suspende a P0 antes de \texttt{turno := 0}}} \\
4 & P1 & \texttt{c[1] := 0} & 0 \; 0 \; 1 & \\
5 & P1 & \texttt{turno <> 1}: falso & 0 \; 0 \; 1 & \textbf{P1 entra en la SC} \\
6 & P0 & \texttt{turno := 0} & 0 \; 0 \; 0 & \\
7 & P0 & \texttt{turno <> 0}: falso & 0 \; 0 \; 0 & \textbf{P0 entra en la SC} \\
\bottomrule
\end{tabular}
\end{center}
Los dos están en la SC: \textbf{se viola la EM}. Causa: P0 toma el turno sin volver a comprobar la clave de P1, y
P1, al tener ya el turno, entra sin mirar la de P0.
\end{solucion}

\begin{ejercicio}[2: Peterson con las dos primeras asignaciones intercambiadas]
\texttt{turno := 1-i; quiere[i] := true; while quiere[1-i] and turno = 1-i do skip}. ¿Se mantiene la EM?
\end{ejercicio}
\begin{solucion}
No. Inicialmente los dos indicadores son falsos.
\begin{enumerate}[leftmargin=*,noitemsep]
  \item P0: \texttt{turno := 1}.
  \item P1: \texttt{turno := 0}.
  \item P1: \texttt{quiere[1] := true}.
  \item P1 evalúa la guarda: \texttt{quiere[0]} es falso, así que \textbf{P1 entra}.
  \item P0: \texttt{quiere[0] := true}.
  \item P0 evalúa la guarda: \texttt{turno = 0}, distinto de 1, así que la guarda es falsa y \textbf{P0 entra}.
\end{enumerate}
Los dos en la SC. El orden importa porque, con el original, quien escribe el turno el último ya ha publicado su
indicador, y el otro lo verá.
\end{solucion}

\begin{ejercicio}[3: segunda y tercera etapa del refinamiento]
(a) Traza que viole la EM en la 2ª etapa. (b) Traza de interbloqueo en la 3ª.
\end{ejercicio}
\begin{solucion}
(a) $c_1 = c_2 = 1$. P1 evalúa \texttt{while c2 = 0}: falso, sale del bucle. P2 evalúa \texttt{while c1 = 0}:
falso, sale. P1: \texttt{c1 := 0}, entra. P2: \texttt{c2 := 0}, entra. \textbf{Dos procesos en la SC.}\\[3pt]
(b) P1: \texttt{c1 := 0}. P2: \texttt{c2 := 0}. P1 evalúa \texttt{while c2 = 0}: cierto, espera. P2 evalúa
\texttt{while c1 = 0}: cierto, espera. Ninguno cambiará su clave: \textbf{interbloqueo}.
\end{solucion}

\begin{ejercicio}[4: demostrar la exclusión mutua de Peterson para 2 procesos]
\end{ejercicio}
\begin{solucion}
Supongamos P0 y P1 en la SC a la vez. Ninguno ha ejecutado su protocolo de salida, así que
\texttt{solicitado[0] = solicitado[1] = true}. Para salir de su bucle, cada $P_i$ tuvo que ver
$\neg\texttt{solicitado[j]} \lor \texttt{turno} \neq i$; como \texttt{solicitado[j]} es cierto desde antes de que
$P_j$ escribiera el turno, P0 vio $\texttt{turno} \neq 0$ y P1 vio $\texttt{turno} \neq 1$. Sea $P_k$ el último
que escribió \texttt{turno}: escribió \texttt{turno := k} y nadie lo cambió después, así que $P_k$ solo pudo
leer $\texttt{turno} = k$ y no pudo salir de su bucle. Contradicción: $N_{SC} \le 1$.
\end{solucion}

\begin{ejercicio}[5: inanición en el algoritmo de Dijkstra]
Explicar por qué el algoritmo de Dijkstra para $N$ procesos puede producir inanición.
\end{ejercicio}
\begin{solucion}
$P_i$ solo puede apropiarse del turno cuando ve pasivo al que lo tiene. Con tres procesos y \texttt{turno = 1}:
P1 entra, sale (\texttt{c[1] := pasivo}), y antes de que $P_3$ consulte \texttt{c[turno]}, P2 lo hace, ve a P1
pasivo y toma el turno. El patrón se repite con P1 y P2 alternándose; el planificador siempre deja a P3 leer
cuando el poseedor del turno no está pasivo. P3 nunca entra: hay \textbf{inanición}, aunque siempre entra alguien
(\textbf{sin interbloqueo}). Knuth lo corrige con el turno circular.
\end{solucion}

\begin{ejercicio}[6: escenario de Ricart-Agrawala]
P1, P2 y P3 piden la SC a la vez, los tres con \texttt{ns = 1}. ¿En qué orden entran? ¿Cuántos mensajes hay en
total?
\end{ejercicio}
\begin{solucion}
Con números iguales desempata el identificador (menor id, más prioridad). Al recibir \texttt{(pet, 1, j)}:
\begin{itemize}[leftmargin=*,noitemsep]
  \item P1 tiene prioridad frente a P2 y P3: \textbf{retrasa} ambas respuestas.
  \item P2 contesta a P1 y retrasa a P3. \quad P3 contesta a P1 y a P2.
\end{itemize}
P1 recibe 2 respuestas y \textbf{entra}. Al salir responde a P2 y P3. P2 completa sus 2 respuestas (de P3 y de P1)
y \textbf{entra}; al salir responde a P3, que entra el último. Orden: \textbf{P1, P2, P3}. Mensajes:
$3 \cdot 2(n-1) = 3 \cdot 4 = \mathbf{12}$. Si P1 vuelve a pedir al salir, su \texttt{ns} será
$\texttt{mns}+1 = 2$ y quedará detrás de las peticiones pendientes.
\end{solucion}

\begin{ejercicio}[7: escenario de Suzuki-Kasami]
$n = 3$. P1 tiene el testigo y no está en la SC; P2 pide la SC por primera vez. ¿Cuántos mensajes hacen falta?
\end{ejercicio}
\begin{solucion}
P2: \texttt{ns := 1}, difunde \texttt{(pet, 1, 2)} a P1 y P3 (2 mensajes). P1 actualiza
\texttt{peticion[2] := 1}; como \texttt{peticion[2] = 1 > token[2] = 0} y no está en la SC, envía el testigo a P2 y
pone \texttt{token\_presente := false} (1 mensaje). Total $\mathbf{3 = n}$. P2 entra y, al salir,
\texttt{token[2] := 1}.
\end{solucion}

\begin{ejercicio}[8: exclusión mutua con \texttt{swap} para $N$ procesos]
\end{ejercicio}
\begin{solucion}
\begin{lstlisting}
-- compartida: ocupado := false ;  local: llave
llave := true;
repeat swap(ocupado, llave) until llave = false;
<<seccion critica>>
ocupado := false;
\end{lstlisting}
\textbf{EM}: tras el primer \texttt{swap} con éxito, \texttt{ocupado} vale cierto; cualquier otro \texttt{swap}
devuelve cierto hasta que el dueño pone \texttt{ocupado := false}. Solo uno recibe falso.
\textbf{Progreso}: si el cerrojo está libre, el primer \texttt{swap} lo toma. \textbf{Inanición}: posible; la
atomicidad no decide quién gana la siguiente carrera. Coste: espera activa.
\end{solucion}

\begin{ejercicio}[9: capacidad $k$ con cerrojo (problema 59)]
\end{ejercicio}
\begin{solucion}
\begin{lstlisting}
entrar():  ok := false;
           repeat
             lock();
             if dentro < k then begin dentro := dentro + 1; ok := true; end;
             unlock();
           until ok;
salir():   lock(); dentro := dentro - 1; unlock();
\end{lstlisting}
\textbf{Invariante} $0 \le \texttt{dentro} \le k$: solo se incrementa bajo el cerrojo tras comprobar
$\texttt{dentro} < k$ con el mismo cerrojo, y solo decrementa quien incrementó antes. \textbf{No se espera con el
cerrojo cogido}: \texttt{salir} lo necesita y nadie podría liberar plazas (interbloqueo). \textbf{Inanición}:
posible, un proceso puede llegar siempre tarde a la plaza libre.
\end{solucion}

% ==============================================================================
\clearpage
\section{Tabla resumen para el último repaso}
\begin{center}\small
\setlength{\tabcolsep}{4pt}
\begin{tabular}{lcccc p{4.2cm}}
\toprule
\textbf{Algoritmo} & \textbf{N} & \textbf{EM} & \textbf{Sin interbl.} & \textbf{Sin inanición} & \textbf{Clave} \\ \midrule
1ª etapa (turno) & 2 & \si & \no & -- & viola 3ª condición \\
2ª etapa & 2 & \no & -- & -- & comprobar y luego asignar \\
3ª etapa & 2 & \si & \no & \no & interbloqueo \\
4ª etapa (cortesía) & 2 & \si & \no{} (livelock) & \no & ceden a la vez \\
Dekker & 2 & \si & \si & \si & claves + turno \\
Hyman & 2 & \no & -- & -- & toma el turno sin recomprobar \\
Peterson & 2 & \si & \si & \si & espera el último en escribir turno; cota 1 \\
Dijkstra & $N$ & \si & \si & \no & dos barreras \\
Knuth & $N$ & \si & \si & \si & turno circular; no FIFO \\
Peterson & $N$ & \si & \si & \si & etapas; $n(n-1)/2$ turnos \\
Lamport (panadería) & $N$ & \si & \si & \si & FIFO; números sin límite \\
Tickets & 2 & \si & \si & \si & \texttt{n[i]:=1} imprescindible \\
Anillo \texttt{s[i]} & $N$ & \si & \no & -- & alternancia forzosa \\
\texttt{swap} / TAS & $N$ & \si & \si & \no & atómico, no equitativo \\
Ricart-Agrawala & $N$ & \si & \si & \si & $2(n-1)$ mensajes \\
Suzuki-Kasami & $N$ & \si & \si & \si & $n$ mensajes (0 si ya lo tiene) \\
Misra & anillo & -- & -- & -- & regenera el testigo; ping + pong = 0 \\
\bottomrule
\end{tabular}
\end{center}
\textit{«--»: no aplica o no se estudia esa propiedad para ese algoritmo.}

\end{document}
